% from main.tex
%%%%%%%%%%  scan idx 559  |  folio 546  %%%%%%%%%%
appelle équivalence numérique la relation d'équivalence définie par le produit d'intersection \(\langle\, ,\,\rangle\) :

\textbf{Définition 4.5.1.} On dit que deux éléments \(x\) et \(y\) de \(\Gr^i(X)\) \((i\in\mathbb Z)\) (resp. de \(\Gr_i(X)\)) sont numériquement équivalents, s'ils définissent le même homomorphisme \(\Gr_i(X)\longrightarrow\mathbb Z\) (resp. \(\Gr^i(X)\longrightarrow\mathbb Z\)).
On dit que deux \(\mathcal O_X\)-modules inversibles \(L_1\) et \(L_2\) sont numériquement équivalents, si leurs premières classes de Chern \(x,y\in\Gr^1(X)\) sont numériquement équivalentes.

\textbf{Proposition 4.5.2.} Un \(\mathcal O_X\)-module inversible \(L\) est numériquement équivalent à zéro si et seulement si
\begin{equation}\tag{4.5.3}
 \chi(L^n\otimes_{\mathcal O_X}\mathcal O_Z)=\chi(\mathcal O_Z)
\end{equation}
pour tout \(n\in\mathbb Z\) et pour toute courbe intègre fermée \(Z\subset X\).

Cela résulte de 1.1.2 et 4.3.1.

On désigne par \(\Pic^N(X)\) le sous-groupe de \(\Pic(X)\) composé des classes de \(\mathcal O_X\)-modules inversibles numériquement équivalents à zéro. De 2.3.4 on tire

\textbf{Corollaire 4.5.3.} \(\Pic^\tau(X)\subset \Pic^N(X)\).

On va démontrer dans Exp XIII que l'on a même \(\Pic^N(X)=\Pic^\tau(X)\).

\section*{5. L'isomorphisme \(\Pic(X)\cong\Gr^1(X)\)}

\subsection*{5.1}
Soit \(X\) un topos annelé quelconque. Alors on peut définir l'anneau \(K^\circ(X)\) et la filtration \((\Filt^i(X))_{i\in\mathbb Z}\) exactement comme pour les schémas. De même on définit le groupe \(\Pic(X)\) comme le groupe des classes d'isomorphie des \(\mathcal O_X\)-modules inversibles. On désigne par
\begin{equation}\tag{5.1.1}
